\documentclass[../main.tex]{subfiles}
\begin{document}

\textit{Chương này khảo sát các công cụ hiện có, xác định tác nhân và nêu các yêu cầu chức năng,
phi chức năng. Danh sách đầy đủ (89 yêu cầu chức năng và 12 yêu cầu phi chức năng; tổng số theo
nhóm dưới đây bằng 89) nằm trong tài liệu D1; ở đây chỉ tóm tắt các yêu cầu quan trọng nhất.}

\section{Khảo sát các công cụ hiện có}
\begin{table}[H]
\centering
\caption{So sánh với các công cụ phổ biến}\label{tab:survey}
\renewcommand{\arraystretch}{1.25}
\begin{tabular}{|p{3.4cm}|p{1.9cm}|p{1.9cm}|p{2.3cm}|p{2.8cm}|}
\hline
\textbf{Tiêu chí} & \textbf{Trello} & \textbf{Jira} & \textbf{Bảng tính} & \textbf{Hệ thống của nhóm} \\ \hline
Bảng Kanban & Có & Có & Thủ công & Có \\ \hline
Báo cáo Scrum (burndown, velocity) & Tiện ích thêm & Có & Không & Có \\ \hline
Phụ thuộc task, đường găng & Không & Một phần & Không & Có (DAG) \\ \hline
Điều khiển hoàn toàn bằng bàn phím & Hạn chế & Hạn chế & Có & Có \\ \hline
Dùng khi ngoại tuyến & Không & Không & Có & Dự án cá nhân \\ \hline
AI tóm tắt, biên bản, giao việc & Trả phí & Trả phí & Không & Có \\ \hline
Độ khó cho sinh viên & Thấp & Cao & Thấp & Thấp \\ \hline
\end{tabular}
\end{table}

Khảo sát cho thấy khoảng trống mà dự án nhắm tới: một công cụ nhẹ như Trello nhưng vẫn có báo cáo
Scrum, phân tích phụ thuộc và hỗ trợ AI, đồng thời dùng được khi không có mạng.

\section{Tác nhân}
\begin{table}[H]
\centering
\caption{Các tác nhân của hệ thống}\label{tab:actors}
\renewcommand{\arraystretch}{1.25}
\begin{tabular}{|p{3.2cm}|p{10.5cm}|}
\hline
\textbf{Tác nhân} & \textbf{Mô tả} \\ \hline
Khách & Chưa đăng nhập. Xem trang giới thiệu (tính năng, video demo, bảng giá, hỏi đáp) và có thể đăng ký hoặc đăng nhập. \\ \hline
Người dùng & Mọi tài khoản đã đăng nhập. Bản thân tài khoản \emph{không} có quyền gì: chỉ có thể tạo dự án, trả lời lời mời và sửa hồ sơ của mình. \\ \hline
Quản lý dự án (PM) & Người dùng giữ vai trò PM \emph{trong một dự án cụ thể}. Quản lý thành viên và vai trò, lập kế hoạch và điều hành sprint. Người tạo dự án trở thành PM của dự án đó. \\ \hline
Thành viên & Người dùng giữ vai trò Thành viên trong dự án. Tạo và cập nhật task, bình luận, dùng công cụ AI; không được đổi vai trò hay sprint. \\ \hline
Dịch vụ AI & LLM bên ngoài (Groq) hoặc mô hình cục bộ (Ollama) trả lời các prompt có cấu trúc. \\ \hline
\end{tabular}
\end{table}

Vai trò được đặt theo từng dự án một cách có chủ đích: cùng một người có thể là PM ở dự án này và là
Thành viên ở dự án khác. Khi đăng ký, hệ thống không bao giờ hỏi vai trò.

\begin{figure}[H]
\centering
\begin{tikzpicture}[
  actor/.style={draw, rounded corners, minimum width=2.4cm, minimum height=0.8cm, fill=gray!10, font=\small},
  uc/.style={draw, ellipse, minimum width=5.6cm, minimum height=0.8cm, inner xsep=0pt, font=\footnotesize, align=center},
  >=Stealth]
\node[actor] (pm) at (0,1.9) {Quản lý dự án};
\node[actor] (mem) at (0,-2.4) {Thành viên};
\node[actor] (ai) at (12.0,-3.2) {Dịch vụ AI};
\node[uc] (u1) at (6.4,4.50) {Quản lý thành viên, vai trò};
\node[uc] (u2) at (6.4,3.00) {Lập, bắt đầu, kết thúc sprint};
\node[uc] (u3) at (6.4,1.50) {Tạo / sửa / chuyển task};
\node[uc] (u4) at (6.4,0.00) {Bình luận, @mention, minh chứng};
\node[uc] (u5) at (6.4,-1.50) {Xem burndown, velocity, DAG};
\node[uc] (u6) at (6.4,-3.00) {AI tóm tắt, biên bản, giao việc};
\node[uc] (u7) at (6.4,-4.50) {AI chia nhỏ, retrospective};
\draw[->] (pm.east) -- (u1.west); \draw[->] (pm.east) -- (u2.west);
\foreach \u in {u3,u4,u5,u6,u7} { \draw[->] (pm.east) -- (\u.west); \draw[->] (mem.east) -- (\u.west); }
\draw[->,dashed] (u6.east) -- (ai.west); \draw[->,dashed] (u7.east) -- (ai.west);
\end{tikzpicture}
\caption{Các ca sử dụng chính (PM làm được mọi việc của Thành viên)}\label{fig:usecase}
\end{figure}

\section{Phân tích nghiệp vụ}
\subsection{Quy trình hiện tại và các vấn đề}
Một nhóm sinh viên thường phối hợp qua nhóm chat và một bảng tính dùng chung. Quy trình này có bốn
vấn đề lặp lại:
\begin{itemize}
  \item \textbf{Không có một nguồn sự thật duy nhất.} Việc được thống nhất trong chat rồi chép tay
        sang bảng tính; hai nơi lệch nhau và không ai biết bản nào mới nhất.
  \item \textbf{Vướng mắc không ai thấy.} Thành viên bị kẹt nhắn một tin rồi tin bị trôi; trưởng
        nhóm chỉ biết khi đến hạn.
  \item \textbf{Không có thứ tự.} Các việc phụ thuộc nhau được liệt kê ngang hàng, nên có người bắt
        đầu việc khi việc tiên quyết chưa xong.
  \item \textbf{Báo cáo thủ công.} Báo cáo tuần và biên bản họp viết theo trí nhớ; việc giao nhiệm vụ
        theo thói quen chứ không theo kỹ năng hay khối lượng.
\end{itemize}

\subsection{Quy trình mục tiêu}
Hệ thống hỗ trợ chu trình theo kiểu Scrum ở Hình~\ref{fig:bizflow}. Mỗi bước là một hoạt động nghiệp
vụ do một vai trò thực hiện, và tạo ra dữ liệu cho bước sau.
\begin{figure}[H]
\centering
\begin{tikzpicture}[
  s/.style={draw, rounded corners, fill=blue!6, minimum width=2.7cm, minimum height=1.15cm, align=center, font=\footnotesize},
  >=Stealth]
\node[s] (a) at (0,0)    {1. Tạo dự án,\\mời thành viên\\\scriptsize(PM)};
\node[s] (b) at (3.6,0)  {2. Lập backlog,\\đặt phụ thuộc\\\scriptsize(PM, Thành viên)};
\node[s] (c) at (7.2,0)  {3. Lập kế hoạch\\sprint, giao việc\\\scriptsize(PM + AI)};
\node[s] (d) at (10.8,0) {4. Bắt đầu sprint\\(ghi nhận\\cam kết)\\\scriptsize(PM)};
\node[s] (e) at (10.8,-2.5) {5. Làm việc hằng\\ngày: trạng thái,\\bình luận, minh chứng\\\scriptsize(Thành viên)};
\node[s] (f) at (7.2,-2.5) {6. Theo dõi:\\burndown, bị chặn,\\quá hạn\\\scriptsize(PM)};
\node[s] (g) at (3.6,-2.5) {7. Kết thúc sprint,\\chuyển việc dở\\\scriptsize(PM)};
\node[s] (h) at (0,-2.5) {8. Retrospective,\\báo cáo tuần\\\scriptsize(PM + AI)};
\draw[->] (a)--(b); \draw[->] (b)--(c); \draw[->] (c)--(d); \draw[->] (d)--(e);
\draw[->] (e)--(f); \draw[->] (f)--(g); \draw[->] (g)--(h);
\draw[->, dashed] (h.north) to[bend left=12] node[left, font=\scriptsize]{sprint sau} (b.south west);
\end{tikzpicture}
\caption{Quy trình nghiệp vụ mục tiêu của một nhóm dự án}\label{fig:bizflow}
\end{figure}

\subsection{Vòng đời của task}
Một task đi qua bốn trạng thái. Thứ tự cố định và giống nhau ở client và server:
\begin{figure}[H]
\centering
\begin{tikzpicture}[st/.style={draw, rounded corners, minimum width=2.6cm, minimum height=0.85cm, font=\small}, >=Stealth]
\node[st, fill=gray!10] (t) at (0,0) {TODO};
\node[st, fill=blue!8] (i) at (3.8,0) {IN\_PROGRESS};
\node[st, fill=red!10] (b) at (7.6,0) {BLOCKED};
\node[st, fill=green!12] (d) at (11.4,0) {DONE};
\draw[->] (t)--(i); \draw[->] (i)--(b); \draw[->] (b)--(d);
\draw[->] (d.south) to[bend left=18] node[below, font=\scriptsize]{mở lại} (t.south);
\end{tikzpicture}
\caption{Vòng trạng thái task (phím Space chuyển một bước)}\label{fig:taskcycle}
\end{figure}

\subsection{Quy tắc nghiệp vụ}
\begin{longtable}{|p{1.3cm}|p{12.4cm}|}
\caption{Các quy tắc nghiệp vụ chính}\label{tab:rules}\\
\hline
\textbf{Mã} & \textbf{Quy tắc} \\ \hline
\endfirsthead
\hline
\textbf{Mã} & \textbf{Quy tắc} \\ \hline
\endhead
BR-01 & Quyền hạn thuộc về vai trò \emph{trong một dự án}; tài khoản tự thân không có quyền gì. Người tạo dự án là PM của dự án đó. \\ \hline
BR-02 & Người được mời chỉ có quyền truy cập sau khi chấp nhận; lời mời đang chờ không cấp quyền gì. \\ \hline
BR-03 & Chỉ thành viên thấy dự án; với người khác, dự án như không tồn tại (404). \\ \hline
BR-04 & Chỉ PM quản lý thành viên, vai trò và sprint. \\ \hline
BR-05 & Mỗi dự án có tối đa một sprint đang chạy; cam kết (story point) được chốt khi sprint bắt đầu. \\ \hline
BR-06 & Khi kết thúc sprint, task chưa xong chuyển sang sprint sau hoặc backlog; xoá sprint không bao giờ xoá công việc. \\ \hline
BR-07 & Task bị chặn nên ghi rõ lý do; task đã xong không bao giờ hiển thị là quá hạn. \\ \hline
BR-08 & Mọi người được giao phải là thành viên đã chấp nhận của dự án; phụ thuộc chỉ trong cùng một dự án. \\ \hline
BR-09 & Thứ tự trên bảng: trễ hạn trước, rồi hạn gần hơn, rồi đến độ ưu tiên. \\ \hline
BR-10 & Dự án cá nhân sửa được khi mất mạng; dự án nhóm chỉ xem khi mất mạng để tránh ghi đè lẫn nhau. \\ \hline
BR-11 & Không ai nhận thông báo về hành động của chính mình. \\ \hline
\end{longtable}

\section{Yêu cầu chức năng}
\begin{longtable}{|p{2.3cm}|p{9.2cm}|p{2.0cm}|}
\caption{Các nhóm yêu cầu chức năng}\label{tab:fr}\\
\hline
\textbf{Nhóm} & \textbf{Yêu cầu tiêu biểu} & \textbf{Số lượng} \\ \hline
\endfirsthead
\hline
\textbf{Nhóm} & \textbf{Yêu cầu tiêu biểu} & \textbf{Số lượng} \\ \hline
\endhead
FR-AUTH & Đăng ký bằng email và mật khẩu (không chọn vai trò); đăng nhập; phiên JWT; vai trò theo dự án; người ngoài nhận 404 với dự án họ không được xem. & 10 \\ \hline
FR-PROJ & Bảng điều khiển cá nhân; tạo, xoá dự án; mời qua email (chờ chấp nhận); đổi vai trò; trang hồ sơ có ảnh đại diện. & 11 \\ \hline
FR-DOM & Task có tiêu đề, mô tả, trạng thái, ưu tiên, độ phức tạp, hạn chót, lý do bị chặn, nhiều người thực hiện, phụ thuộc, tiêu chí chấp nhận; bình luận có @mention, trả lời một cấp, đính kèm; thông báo. & 10 \\ \hline
FR-INT & Kanban và dạng bảng; mọi thao tác bằng phím ($h/j/k/l$, Space đổi trạng thái, $n$ tạo mới, Enter xem chi tiết, $/$ tìm kiếm, 1--4 đổi ưu tiên); sắp xếp theo hạn chót trước. & 14 \\ \hline
FR-GRAPH & Thứ tự topo, chịu được chu trình, đường găng hiển thị trên bảng và màn hình DAG. & 5 \\ \hline
FR-PERS & IndexedDB trước, server sau; dự án cá nhân sửa được khi ngoại tuyến, dự án chung chỉ đọc. & 6 \\ \hline
FR-SCRUM & Vòng đời sprint (đã lên kế hoạch, đang chạy, đã xong); mỗi dự án tối đa một sprint đang chạy; ghi nhận cam kết khi bắt đầu; chuyển việc chưa xong khi kết thúc; lập kế hoạch backlog; burndown; velocity; retrospective AI. & 8 \\ \hline
FR-HIST & Mọi thay đổi của task được ghi kèm người sửa và thời điểm; hiển thị trong màn hình chi tiết. & 2 \\ \hline
FR-AI & Tóm tắt tuần; biên bản họp; gợi ý giao việc; chia nhỏ mục tiêu; phân tích git diff; thống kê khối lượng và task trễ hạn; dự phòng ba tầng. & 8 \\ \hline
FR-NAV, FR-I18N, FR-MKT & Trang giới thiệu với nút đăng nhập góc trên bên phải; chuyển Anh/Việt cho toàn bộ giao diện; video demo quay từ ứng dụng thật. & 15 \\ \hline
\end{longtable}

\subsection{Quy trình Scrum}
Các yêu cầu Scrum định nghĩa một máy trạng thái nhỏ cho sprint:
\begin{figure}[H]
\centering
\begin{tikzpicture}[st/.style={draw, rounded corners, minimum width=2.6cm, minimum height=0.9cm, font=\small}, >=Stealth, node distance=3.3cm]
\node[st] (p) {PLANNED};
\node[st, right=of p] (a) {ACTIVE};
\node[st, right=of a] (c) {COMPLETED};
\draw[->] (p) -- node[above, font=\scriptsize]{bắt đầu (PM)} node[below, font=\scriptsize]{lưu cam kết} (a);
\draw[->] (a) -- node[above, font=\scriptsize]{kết thúc (PM)} node[below, font=\scriptsize, align=center]{việc chưa xong $\rightarrow$\\sprint sau / backlog} (c);
\end{tikzpicture}
\caption{Vòng đời sprint}\label{fig:sprint-states}
\end{figure}
Chỉ PM được thay đổi sprint; mỗi dự án tối đa một sprint đang chạy; sprint đã xong không nhận thêm
việc; xoá sprint sẽ trả task về backlog chứ không bao giờ xoá công việc.

\section{Yêu cầu phi chức năng}
\begin{table}[H]
\centering
\caption{Các yêu cầu phi chức năng chính}\label{tab:nfr}
\renewcommand{\arraystretch}{1.25}
\begin{tabular}{|p{2.9cm}|p{10.8cm}|}
\hline
\textbf{Thuộc tính} & \textbf{Yêu cầu} \\ \hline
Độ phản hồi & Thay đổi cục bộ hiển thị ngay; lời gọi mạng không bao giờ chặn việc hiển thị. \\ \hline
Tính sẵn sàng của AI & Không endpoint AI nào trả lỗi vì mô hình không sẵn sàng; thay vào đó trả câu trả lời tất định và ghi log tình trạng suy giảm. \\ \hline
Bảo mật & Server kiểm tra quyền ở mọi route theo dự án; ngoài môi trường phát triển, dịch vụ từ chối khởi động nếu khoá bí mật yếu, còn tài khoản demo, CORS dạng wildcard hoặc lược đồ CSDL chưa cập nhật. \\ \hline
Toàn vẹn dữ liệu & Chỉ thay đổi lược đồ qua migration có phiên bản; bắt buộc khoá ngoại; UTF-8 (utf8mb4) xuyên suốt để tiếng Việt không bị cắt. \\ \hline
Đa ngôn ngữ & Mọi chuỗi hiển thị đều có bản tiếng Anh và tiếng Việt; câu trả lời của AI theo ngôn ngữ giao diện. \\ \hline
Khả năng kiểm thử & Bộ test backend và frontend chạy được trên bản checkout sạch và chạy trong CI. \\ \hline
\end{tabular}
\end{table}

\end{document}
